
\subsection*{6.4. Démonstration de 6.2.2.}
Il suffit de montrer que, quels que soient @@M0@@, @@M1@@ et @@M2@@, on peut trouver un «~voisinage étale~» @@M3@@ de @@M4@@ tel que la restriction de @@M5@@ sur @@M6@@ soit un produit de commutateurs de @@M7@@. La question étant locale, on peut supposer @@M8@@ affine d'anneau @@M9@@, et @@M10@@ constant égal à @@M11@@ sur @@M12@@. Soient @@M13@@ l'idéal premier de @@M14@@ et @@M15@@ la matrice de @@M16@@. Pour tout @@M17@@ et @@M18@@ ; @@M19@@, nous désignerons par @@M20@@ la matrice @@M21@@ @@M22@@ telle que @@M23@@ et @@M24@@ pour @@M25@@ @@M26@@ symbole de Kronecker). Pour toute matrice @@M27@@ on obtient @@M28@@ en ajoutant à la @@M29@@-ème colonne de @@M30@@ @@M31@@ fois la @@M32@@-ème. De même on obtient la matrice @@M33@@ en ajoutant à la @@M34@@-ème ligne de @@M35@@ @@M36@@ fois la @@M37@@-ème. En particulier on a
@@M38@@
et on voit aussitôt que l'on a
@@M39@@
pour tout @@M40@@, de sorte qu'on trouve une représentation
@@M41@@
Sauf mention expresse du contraire, on va supposer dans ce qui suit le @@M42@@ intervenant dans @@M43@@ @@M44@@ inversible.
La démonstration se fait alors en deux pas. On va montrer d'abord que @@M45@@ est un produit de matrices @@M46@@ dans un voisinage zariskien suffisamment petit de @@M47@@. Puis on va montrer que les matrices @@M48@@ sont des produits de commutateurs de @@M49@@ dans un voisinage étale suffisamment petit de @@M50@@.

